\section{Experiments}
\label{sec:exp}
In this section, we probe the empirical performance of MES and add-MES on a variety of tasks. Here, MES-G denotes MES with $y_*$ sampled from the approximate Gumbel distribution, and MES-R denotes MES with $y_*$ computed by maximizing a sampled function represented by random features. Following~\cite{hennig2012,hernandez2014predictive}, we adopt the zero mean function and non-isotropic squared exponential kernel as the prior for the GP. We compare to methods from the entropy search family, i.e., ES and PES, and to other popular Bayesian optimization methods including GP-UCB (denoted by UCB), PI, EI and EST. The parameter for GP-UCB was set according to Theorem 2 in~\cite{srinivas2009gaussian};
%
the parameter for PI was set to be the observation noise $\sigma$. For the functions with unknown GP hyper-parameters, every 10 iterations, we learn the GP hyper-parameters using the same approach as was used by PES~\cite{hernandez2014predictive}. For the high dimensional tasks, we follow~\cite{kandasamy2015high} and sample the additive structure/GP parameters with the highest data likelihood when they are unknown. We evaluate performance according to the simple regret and inference regret as defined in Section~\ref{ssec:eval}. We used the open source Matlab implementation of PES, ES and EST~\cite{hennig2012,hernandez2014predictive,wang2016est}.
 Our Matlab code and test functions are available at \url{https://github.com/zi-w/Max-value-Entropy-Search/}.
\subsection{Synthetic Functions}
\begin{figure}
\centering
\includegraphics[width=.5\textwidth]{figs/3d_synth_new4}
\caption{(a) Inference regret; (b) simple regret. MES methods are much less sensitive to the number of maxima $y_*$ sampled for the acquisition function (1, 10 or 100) than PES is to the number of argmaxes $x_*$. }
\label{fig:synth1}
\end{figure}

\begin{table}[t]
\caption{The runtime of selecting the next input. PES 100 is significantly slower than other methods. MES-G's runtime is comparable to the fastest method EI while it performs better in terms of simple and inference regrets.}
\label{tb:botiming}
\vskip -0.2in
\begin{center}
\begin{small}
\begin{sc}
\begin{tabular}{lc|lc}%
\hline
\abovespace\belowspace
Method & Time (s) & Method & Time (s) \\
\hline
\abovespace
UCB    & {\color{red}$0.08\pm0.05$}  & PES 1& $0.20\pm 0.06$ \\
PI & {\color{red}$0.10\pm0.02 $}        & MES-R 100& $5.85\pm 0.86 $\\
EI    & {\color{red}$0.07\pm 0.03$}             & MES-R 10 & $0.67\pm 0.11$ \\
EST    & $0.15\pm0.02$    & MES-R 1&  $0.13\pm 0.03$       \\
ES    & $8.07\pm3.02$      &  MES-G 100 & $0.12\pm 0.02$\\
PES 100   & $15.24\pm4.44$    & MES-G 10 & {\color{red}$0.09\pm0.02 $} \\
PES 10    & $1.61\pm 0.50$      & MES-G 1&    {\color{red}$0.09\pm0.03$}     \\
\hline
\end{tabular}
\end{sc}
\end{small}
\end{center}
\vskip -0.2in
\end{table}

We begin with a comparison on synthetic functions sampled from a 3-dimensional GP, to probe our conjecture that MES is much more robust to the number of $y_*$ sampled to estimate the acquisition function than PES is to the number of $x_*$ samples.
%
For PES, we sample 100 (PES 100), 10 (PES 10) and 1 (PES 1) argmaxes for the acquisition function. Similarly, we sample 100, 10, 1 $y_*$ values for MES-R and MES-G. We average the results on 100 functions sampled from the same Gaussian kernel with scale parameter $5.0$ and bandwidth parameter $0.0625$, and
%
observation noise $\mathcal{N}(0,0.01^2)$.

Figure~\ref{fig:synth1} shows the simple and inference regrets. 
For both regret measures, PES is very sensitive to the the number of $x_*$ sampled for the acquisition function: 100 samples lead to much better results than 10 or 1. In contrast, both MES-G and MES-R perform competitively even with 1 or 10 samples.
%
Overall, MES-G is slightly better than MES-R, and both MES methods performed better than other ES methods. 
MES methods performed better than all other methods with respect to simple regret. For inference regret, MES methods performed similarly to EST, and much better than all other methods including PES and ES. %

In Table~\ref{tb:botiming}, we show the runtime of selecting the next input per iteration\footnote{All the timing experiments were run exclusively on an Intel(R) Xeon(R) CPU E5-2680 v4 @ 2.40GHz. The function evaluation time is excluded.} using GP-UCB, PI, EI, EST, ES, PES, MES-R and MES-G on the synthetic data with fixed GP hyper-parameters. For PES and MES-R, every $x_*$ or $y_*$ requires running an optimization sub-procedure, so their running time grows noticeably with the number of samples. MES-G avoids this optimization, and competes with the fastest methods EI and UCB.

In the following experiments, we set the number of $x_*$ sampled for PES to be 200, and the number of $y_*$ sampled for MES-R and MES-G to be 100 unless otherwise mentioned. 
%



\subsection{Optimization Test Functions}
\hide{
\begin{figure*}
\centering
\includegraphics[width=0.85\textwidth]{figs/testfuncs}
\caption{(a) 2-d eggholder function; (b) 10-d Shekel function; (c) 10-d Michalewicz function. PES achieves lower regret on the 2-d function while MES-G performed better than other methods on the two 10-d optimization test functions.}
\label{fig:test}
\end{figure*}
}
We test on three challenging optimization test functions: the 2-dimensional eggholder function, the 10-dimensional Shekel function and the 10-dimensional Michalewicz function. All of these functions have many local optima. We randomly sample 1000 points to learn a good GP hyper-parameter setting, and then run the BO methods with the same hyper-parameters. The first observation is the same for all methods.
%
We repeat the experiments 10 times. The averaged simple regret is shown in the appendix%
, and the inference regret is shown in Table~\ref{tb:test}. On the 2-d eggholder function, PES was able to achieve better function values faster than all other methods, which verified the good performance of PES when sufficiently many $x_*$ are sampled. However, for higher-dimensional test functions, the 10-d Shekel and 10-d Michalewicz function, MES methods performed much better than PES and ES, and MES-G performed better than all other methods. 

\begin{table}[t]
\caption{Inference regret $R_T$ for optimizing the eggholder function, Shekel function, and Michalewicz function.}
\label{tb:test}
\vskip -0.5in
\begin{center}
\begin{small}
\begin{sc}
\begin{tabular}{llll}%
\hline
\abovespace\belowspace
Method & Eggholder & Shekel & Michalewicz \\
\hline
\abovespace
UCB & $141.00 \pm 70.96$ & $9.40 \pm 0.26$ & $6.07 \pm 0.53$   \\
PI & $52.04 \pm 39.03$ & $6.64 \pm 2.00$ & $4.97 \pm 0.39$   \\
EI & $71.18 \pm 59.18$ & $6.63 \pm 0.87$ & $4.80 \pm 0.60$   \\
EST & $55.84 \pm 24.85$ & $5.57 \pm 2.56$ & $5.33 \pm 0.46$   \\
ES & $48.85 \pm 29.11$ & $6.43 \pm 2.73$ & $5.11 \pm 0.73$   \\
PES & {\color{red}$37.94 \pm 26.05$} & $8.73 \pm 0.67$ & $5.17 \pm 0.74$   \\
MES-R & $54.47 \pm 37.71$ & $6.17 \pm 1.80$ & $4.97 \pm 0.59$   \\
\belowspace
MES-G & $46.56 \pm 27.05$ & {\color{red}$5.45 \pm 2.07$} & {\color{red}$4.49 \pm 0.51$}   \\
\hline
\end{tabular}
\end{sc}
\end{small}
\end{center}
\vskip -0.2in
\end{table}

\subsection{Tuning Hyper-parameters for Neural Networks}
\begin{figure}
\centering
\includegraphics[width=0.45\textwidth]{figs/nnet_mean}
\caption{Tuning hyper-parameters for training a neural network, (a) Boston housing dataset; (b) breast cancer dataset. MES methods perform better than other methods on (a), while for (b), MES-G, UCB, PES perform similarly and better than others.}
\label{fig:nnet}
\end{figure}


\begin{table}[t]
\caption{Inference regret $R_T$ for tuning neural network hyper-parameters on the Boston housing and breast cancer datasets.}
\label{tb:nnet}
\vskip -0.5in
\begin{center}
\begin{small}
\begin{sc}
\begin{tabular}{lll}%
\hline
\abovespace\belowspace
Method & Boston & Cancer (\%) \\
\hline
\abovespace
UCB & $1.64 \pm 0.43$ & {\color{red}$3.83 \pm 0.01$}  \\
PI & $2.15 \pm 0.99$ & $4.40 \pm 0.01$  \\
EI & $1.99 \pm 1.03$ & $4.40 \pm 0.01$  \\
EST & $1.65 \pm 0.57$ & $3.93 \pm 0.01$  \\
ES & $1.79 \pm 0.61$ & $4.14 \pm 0.00$  \\
PES &  {\color{red}$1.52 \pm 0.32$} &  {\color{red}$3.84 \pm 0.01$}  \\
MES-R & $1.54 \pm 0.56$ & $3.96 \pm 0.01$  \\
\belowspace
MES-G & {\color{red}$1.51 \pm 0.61$} & {\color{red}$3.83 \pm 0.01$}  \\
\hline
\end{tabular}
\end{sc}
\end{small}
\end{center}
\vskip -0.2in
\end{table}

Next, we experiment with Levenberg-Marquardt optimization for training a 1-hidden-layer neural network. The 4 parameters we tune with BO are the number of neurons,  the damping factor $\mu$, the $\mu$-decrease factor, and  the $\mu$-increase factor. We test regression on the Boston housing dataset and classification on the breast cancer dataset~\cite{bache2013uci}. The experiments are repeated 20 times, and the neural network's weight initialization and all other parameters are set to be the same to ensure a fair comparison. Both of the datasets were randomly split into train/validation/test sets. We initialize the observation set to have 10 random function evaluations which were set to be the same across all the methods. The averaged simple regret for the regression L2-loss on the validation set of the Boston housing dataset is shown in Fig.~\ref{fig:nnet}(a), and the classification accuracy on the validation set of the breast cancer dataset is shown in Fig.~\ref{fig:nnet}(b).   For the classification problem on the breast cancer dataset, MES-G, PES and UCB achieved a similar simple regret. On the Boston housing dataset, MES methods achieved a lower simple regret. We also show the inference regrets for both datasets in Table~\ref{tb:nnet}. 


\subsection{Active Learning for Robot Pushing}
\begin{figure*}
\centering
\includegraphics[width=0.95\textwidth]{figs/highsynth}
\caption{Simple regrets for add-GP-UCB and add-MES methods on the synthetic add-GP functions. Both add-MES methods  outperform add-GP-UCB except for add-MES-G on the input dimension $d=100$. Add-MES-G achieves the lowest simple regret when $d$ is relatively low, while for higher $d$ add-MES-R becomes better than add-MES-G.}
\label{fig:highsynth}
\end{figure*}
\begin{figure}
\centering
\includegraphics[width=.45\textwidth]{figs/robotpush2}
\caption{BO for active data selection on two robot pushing tasks for minimizing the distance to a random goal with (a) 3-D actions and (b) 4-D actions. MES methods perform better than other methods on the 3-D function. For the 4-D function, MES methods converge faster to a good regret, while PI achieves lower regret in the very end.}
\label{fig:robot}
\end{figure}
We use BO to do active learning for the pre-image learning problem for pushing~\cite{kaelbling17icra}. The function we optimize takes as input the pushing action of the robot, and outputs the distance of the pushed object to the goal location. We use BO to minimize the function in order to find a good pre-image for pushing the object to the designated goal location. The first function we tested has a 3-dimensional input: robot location $(r_x,r_y)$ and pushing duration $t_r$. We initialize the observation size to be one, the same across all methods. The second function has a 4-dimensional input: robot location and angle $(r_x,r_y, r_{\theta})$, and pushing duration $t_r$. We initialize the observation to be 50 random points and set them the same for all the methods. We select 20 random goal locations for each function to test if BO can learn where to push for these locations. We show the simple regret in Fig.~\ref{fig:robot} and the inference regret in Table~\ref{tb:robot}. MES methods performed on a par with or better than their competitors. 

%

\begin{table}[t]
\caption{Inference regret $R_T$ for action selection in robot pushing.}
\label{tb:robot}
\vskip -0.5in
\begin{center}
\begin{small}
\begin{sc}
\begin{tabular}{lll}%
\hline
\abovespace\belowspace
Method & 3-d action & 4-d action \\
\hline
\abovespace
UCB & $1.10 \pm 0.66$ & $0.56 \pm 0.44$  \\
PI & $2.03 \pm 1.77$ & {\color{red}$0.16 \pm 0.20$}  \\
EI & $1.89 \pm 1.87$ & $0.30 \pm 0.33$  \\
EST & $0.70 \pm 0.90$ & $0.24 \pm 0.17$  \\
ES &  {\color{red}$0.62 \pm 0.59$} & $0.25 \pm 0.20$  \\
PES & $0.81 \pm 1.27$ & $0.38 \pm 0.38$  \\
MES-R & {\color{red}$0.61 \pm 1.23$} & {\color{red}$0.16 \pm 0.10$}  \\
\belowspace
MES-G & {\color{red}$0.61 \pm 1.26$} & $0.24 \pm 0.25$  \\
\hline
\end{tabular}
\end{sc}
\end{small}
\end{center}
\vskip -0.2in
\end{table}

